\scanpage{10}

\uline{Il y a tjrs une forme de polarisation $P$ sur} $M$, i.e. \deleted{en effet, il suffit de} une polarisation $\varphi_0$ de $M_0$ invariante par $J_M$. En effet, si $\psi_0$ est une polarisation de $M_0$, on pose

\[
\begin{aligned}
\varphi_0(x,y)&=\psi_0(x,y)+\psi_0(J_Mx,J_My)\\
\text{donc}\quad\varphi_0(J_Mx,J_My)&=\psi_0(J_Mx,J_My)+\psi_0(-x,-y)=\varphi_0(x,y).
\end{aligned}
\]

Par suite, on a $\check M\simeq\overline M(-i)$ pour tt $M\in\operatorname{ob}\mathcal M_{\mathbb C}^i$ ).

13). On voit que pour une $\mathcal M$ donnée \inserted{sur $\mathbb R$}, on a des implications

\diagram{\begin{tikzcd}[ampersand replacement=\&,row sep=small,column sep=small]
{}\arrow[d,Rightarrow] \& \alpha)\quad\mathcal M\text{ admet une polarisation} \& {}\\
{}\arrow[d,Rightarrow] \& \beta)\quad\forall M\in\operatorname{Ob}\mathcal M_{\mathbb C}^{i},\ \text{on a }\check M\simeq\overline M(-i), \& \inserted{\mathcal M\ \frac12\text{ simple}}\\
{} \& \gamma)\quad\forall M\in\operatorname{Ob}\mathcal M_{\deleted{\unsure{\mathbb R|\beta}}}^{i},\ \text{on a }\check M\simeq M(-i), \& \inserted{\mathcal M\ \frac12\text{ simple}}
\end{tikzcd}}

($\beta\Rightarrow\gamma$ provient du fait que $M_{\mathbb C}\simeq N_{\mathbb C}$ implique $M\simeq N$).

\margin{\uline{Lemme}}

Les implications précédentes \deleted{sont} \unsure{s'illustrent|sont strictes} chez \inserted{au moins}, $\mathcal M=\operatorname{Rep}(G)$, $G$ groupe algébrique \inserted{affine \deleted{\unsure{connexe} \illegible{5}}} \deleted{\unsure{simple}} sur $\mathbb R$. On peut avoir $\beta$) \deleted{\unsure{et non|sans avoir}} sans avoir $\alpha$)

\margin{?}

[$\deleted{\operatorname{SO}(\varphi)}$ \deleted{\unsure{de}} $\frac{\operatorname{SO}(\varphi)}{\text{Centre}}$, $\varphi$ forme à un nb pair de variables, avec un nombre impair de carrés négatifs], et $\gamma$) sans avoir $\beta$) [$\ldotp\ldotp\ldotp\ldotp$\quad ?\quad$\ldotp\ldotp\ldotp\ldotp\ldotp$]

% archive sheet number: 9
